\documentclass[11pt,a4paper]{article}
\usepackage[margin=2.5cm]{geometry}
\usepackage{amsmath,amssymb}
\usepackage[hidelinks]{hyperref}

\title{Updates to the single crystal models of NCrystal\\
  since the elastic paper}
\author{NCrystal developers}
\date{NCrystal 5.0.0}

\newcommand{\Eq}[1]{Eq.~(#1) of~\cite{elastic}}
\newcommand{\Eqs}[1]{Eqs.~(#1) of~\cite{elastic}}

\begin{document}
\maketitle

\noindent This note summarises how the Bragg diffraction models for single
crystals with isotropic Gaussian mosaicity (\texttt{SCBragg}) and layered
single crystals (\texttt{LCBragg}) in NCrystal 5.0.0 differ from their
description in the NCrystal elastic paper~\cite{elastic} (sections 4 and 5
there, valid for NCrystal 2.0--4.4). The physics model is unchanged: the
cross section of a plane is still $\sigma^{hkl} = q_{hkl}\,g_{hkl}$
(\Eqs{8--9}), with the truncated Gaussian mosaicity distribution $W(\delta)$
of \Eq{18} (with truncation angle $\tau$ from \Eq{17}), and the geometrical
factor $g_{hkl}$ of \Eq{24}. What changed is mostly how $g_{hkl}$ and its
integral over crystallite rotations, $g^\text{LC}_{hkl}$ (\Eq{51}), are
evaluated: faster, and in places more precisely. There is also a bug fix.
Notation is that of~\cite{elastic} throughout: $\alpha = \pi/2-\theta_B$ with
$\cos\alpha = \lambda/2d_{hkl}$, $\gamma$ is the angle between $\hat n$ and
$-\hat k_i$, $\delta(t)$ the angle between $\hat n$ and the point $\hat p(t)$
on the Bragg circle (\Eq{23}), $t'$ the upper limit of integration (\Eq{30}),
$\sigma$ the mosaicity, and $\epsilon$ the configured precision
(\texttt{mosprec}). The closed-form approximation of
$g_{hkl}$ (\Eqs{35--36}), and the conditions for using it (\Eqs{37--38}), are
unchanged.

\section{Single crystals}

\paragraph{Shared models.} Models are set up in the frame of the crystal, and
shared between all materials differing only in the orientation of the
crystal. A wrapper process (\texttt{RotatedScatter}) rotates neutron
directions to and from that frame. Results are unchanged.

\paragraph{Bug fix near back-scattering.} The condition for a plane to
contribute, \Eq{25}, is always fulfilled when its right hand side,
$\cos\tau-\cos\alpha\cos\gamma$, is negative. Its squared forms,
\Eqs{26--27}, correctly use $\max(0,\ldots)$ on the right hand side, but with
a strict inequality, so they wrongly reject the case $\sin\gamma = 0$ when it
is negative. That is a neutron exactly along a plane normal, near
back-scattering ($\alpha < \tau$). Such planes are no longer ignored, and
cross sections are continuous there.

\paragraph{Preselection of planes.} The test of \Eq{27} was carried out for
every pair of plane normals and every neutron. Now a cheaper test comes
first: a pair of normals can only contribute if
\begin{equation}
  \big|\,|\cos\gamma| - \cos\alpha\,\big| \le 2\sin(\tau/2),
\end{equation}
since $|\cos\gamma-\cos\alpha| = 2|\sin\frac{\gamma+\alpha}{2}
\sin\frac{\gamma-\alpha}{2}| \le 2\sin(\tau/2)$ when $|\alpha-\gamma| \le
\tau$ (and likewise for $\pi-\gamma$). This is evaluated in single precision
(with a safety margin, so it never rejects a contributing plane) for up to 64
planes at a time, in a loop written for auto-vectorisation by compilers. The
exact treatment then follows only for the few surviving planes. Results are
unchanged.

\section{Numerical evaluation of the geometrical integral}

\paragraph{Gauss-Legendre quadrature.} Where the closed form can not be used,
\Eq{24} was integrated with Romberg integration, using a splined
representation of the density (section 4.4 of~\cite{elastic}). It is now
integrated with Gauss-Legendre quadrature, using the exact density. With
$w = 1-\cos\delta(t)$, \Eq{23} gives
\begin{equation}
  w = w_0 + 2\sin\alpha\sin\gamma\sin^2(t/2), \qquad
  w_0 = 1-\cos\delta_0 = \frac{\sin^2(\gamma-\alpha)}{1+\cos(\gamma-\alpha)},
\end{equation}
where the form of $w_0$ avoids cancellations at small angles. The exponent
of \Eq{24} then follows from
$\delta^2 = \sum_{m\ge1} 2^{m+1}w^m/(m^2\binom{2m}{m})$ (as many terms as
needed for a relative precision of $10^{-17}$ at $w = 1-\cos\tau$), or else
from $\delta^2 = 4\arcsin^2\sqrt{w/2}$. The interval $[0,t']$ is split into
$\lceil\tau/10\sigma\rceil$ segments. The number of points per segment
(8--24) is chosen from a table of the maximum errors observed, as a function
of the truncation in units of $\sigma$ per segment, such that these are 10
times below the accuracy target $\min(\max(0.1\epsilon,10^{-7}),10^{-4})$
(which was also used with Romberg integration).

\paragraph{Table of geometrical integrals.} When the closed form is never
valid (i.e.\ $\tau > \psi$, which happens for large mosaicities, or small
values of $\epsilon$), a two-dimensional table of $g_{hkl}$ replaces most
numerical integrations. It covers $1.2\tau \le \alpha \le \pi/2$ and
$|\gamma-\alpha| < \tau$, with coordinates which are cheap to evaluate
(without transcendental functions):
\begin{equation}
  z = \frac{\sin(\gamma-\alpha)}{\sin\tau}, \qquad
  r = \frac{\tan(\alpha/2)}{\tan(\alpha/2)+\kappa}, \quad \kappa = \tan(0.6\tau),
\end{equation}
both with uniformly spaced nodes. It holds the smooth function
$\ln\!\big(g_{hkl}/(\sqrt{1-z^2}\sqrt{\sin\alpha/\sin\gamma})\big)$, which is
interpolated with $4\times4$-point Lagrange polynomials. The factor
$\sqrt{\sin\alpha/\sin\gamma}$ is the one of \Eq{36}, and $\sqrt{1-z^2}$
handles the square-root behaviour at the truncation edges. When set up, the
table ($64\times256$ nodes or more) is validated against the numerical
integration at cell centres, and refined until the maximum error is 10 times
below the accuracy target (ignoring values below $10^{-3}N\sigma$, which are
negligible), or else not used.

\section{Layered single crystals}

With the definitions of section 5.1 of~\cite{elastic}, \Eq{42} reads
$u \equiv \cos\gamma = A\cos\varphi + B$, with
$A = \sin\theta_k\sin\theta_n$ and $B = \cos\theta_k\cos\theta_n$ (for
planes at $\pi-\theta_n$, the sign of $B$ changes). Changing the variable of
integration from $\varphi$ to $u$, \Eq{43} becomes
\begin{equation}
  g^\text{LC}_{hkl} = \frac{1}{\pi}\int g^\text{SC}_{hkl}(u)\,G(u)\,du,
  \qquad G(u) = \big(A^2-(u-B)^2\big)^{-1/2},
\end{equation}
over $u \in [B-A,B+A]$.

\paragraph{Series expansion.} The integrand is only non-zero in a narrow
window, $|u-c| \le h$, around $c = \cos\alpha$, with
$h = \sin\alpha\sin\tau + \cos\alpha(1-\cos\tau)$. If that window lies inside
$[B-A,B+A]$, $G$ is expanded to 6th order around $c$, giving
$g^\text{LC}_{hkl} \approx \frac{1}{\pi}\sum_{n=0}^6 a_n M_n$. With
$d_1 = B+A-c$ and $d_2 = c-B+A$, the Taylor coefficients of $G$ are
$a_n = G_0\sum_{k=0}^n b_k b_{n-k}d_1^{-k}(-d_2)^{-(n-k)}$, where
$G_0 = (d_1d_2)^{-1/2}$ and $b_k = \binom{2k}{k}/4^k$. The moments
$M_n = \int g^\text{SC}_{hkl}(u)(u-c)^n du$ are exactly
\begin{equation}
  \begin{split}
  M_n &= \sin\alpha\,\mathrm{E}_W\Big[\big\langle(-c\,w +
  \sin\alpha\sin\delta\cos\psi)^n\big\rangle_\psi\Big]\\
  &= \sin\alpha\sum_j \binom{n}{2j}\frac{\binom{2j}{j}}{4^j}(-c)^{n-2j}
  \sin^{2j}\!\alpha\;\mathrm{E}_W\big[w^{n-j}(2-w)^j\big],
  \end{split}
\end{equation}
where $\mathrm{E}_W$ is an average over the mosaicity distribution $W(\delta)$,
$w = 1-\cos\delta$, and $\langle\rangle_\psi$ an average over the azimuthal
angle $\psi$ around the nominal normal. This follows from the symmetry of
pairs of points at angular distance $\alpha$ on the unit sphere. Only the
moments $\mathrm{E}_W[w^m]$, $m \le 6$, are needed, and these are computed
once. Since $|a_n| \le G_0/D^n$ with $D = \min(d_1,d_2)$, and
$|M_n| \le h^{n-6}M_6$ for $n > 6$ (as $g^\text{SC}_{hkl} \ge 0$), the
neglected terms are bounded by $G_0M_6\,x/(D^6(1-x))$, with $x = h/D$. The
expansion is used if this bound is at most $\epsilon/10$ of the result, which
is the usual case except for large mosaicities.

\paragraph{Numerical integration.} Otherwise, \Eq{51} is integrated with the
trapezoidal rule, starting with 17 points and halving the step until two
estimates agree to within $\epsilon$, with the integrand evaluated in batches.
This replaces Romberg integration: the integrand is smooth, and either
vanishes or is symmetric at the end points, so that the trapezoidal rule
converges fast and the extrapolations of Romberg integration do not help.

\section{Precision and performance}

The changes in sections 3 and 4 change cross sections by much less than
$\epsilon$ (typically in the 5th significant digit), and generally make them
more precise. Note that the closed-form approximation of \Eq{36}, which is
unchanged, integrates to up to $4\times10^{-5}$ too little near
back-scattering, whatever the value of $\epsilon$. Typical speedups of cross
section evaluations are 2--7 for single crystals (up to 17 for small
mosaicities and low crystal symmetry), 3--15 for layered single crystals, and
5--9 where the closed form is never valid. The sampling of scattering events
(sections 4.5 and 5.3 of~\cite{elastic}) is unchanged.

\section{Configuration}

Not affecting the models themselves, the orientation of single crystals can
now also be given with lattice directions (\texttt{@crys\_uvw:}), with an
additional rotation around a lab axis (\texttt{rot}), or with
\texttt{dir2=arbitrary}, and the axis $\hat L$ of layered crystals can be given
as a plane normal, a lattice direction, or in the lab frame. See the
documentation of the cfg-string parameters.

\begin{thebibliography}{9}
\bibitem{elastic} T.~Kittelmann and X.-X.~Cai, \emph{Elastic neutron
  scattering models for NCrystal}, Comput.\ Phys.\ Commun.\ \textbf{267}
  (2021) 108082, \href{https://doi.org/10.1016/j.cpc.2021.108082}%
  {doi:10.1016/j.cpc.2021.108082}.
\end{thebibliography}

\end{document}
